\section{化学品のGHS：分類と伝達を合わせる}
\label{theory:safety:ghs}
同じ化学品でも国ごとに警告の方法が異なると、運搬・販売・使用の途中で危険を見落とす。
GHSは危険有害性を共通の基準で分類し、その結果をラベルやSDSで伝える枠組みである。
分類の種類と程度を表すクラス・区分、絵表示、注意喚起語、危険有害性情報と注意書きを対応させる。
国内法の表示事項を追加することと国際的な調和は両立する。
当時のGHS案内と平成24年の指針を参照する。\cite{saf53_ghs_old,saf53_ghs_notice}

分類が伝えるのは危険有害性であり、どの作業でも同じ被害確率という意味ではない。
例えば同じ毒性の液体でも、密閉移送と開放した噴霧では曝露の条件が違う。
容器名と供給者が分かることは、追加情報を問い合わせ、事故時に追跡するためにも必要である。
複数のラベル要素の優先規則を、任意に不都合な危険だけ省略できる規則へ変えない。
新しい知見を得たら情報を更新する。当時のJISの細目と現行の表示義務は、版を分けて確認する。

\section{機械安全：本質的な設計と残る危険の情報}
\label{theory:safety:machine}
信頼性は、指定の条件と期間に要求された機能を果たす能力である。
高い信頼性で意図どおり動く装置でも、目的とする動きが人を傷つける設計なら安全とはいえない。
安全性を調べる際は、使用範囲と寿命段階、危険源、危険へ接近する状況を明確にする。
リスク分析の結果を評価し、必要な低減を繰り返す。\cite{saf53_machine2007}

まず危険を生む形状・運動・エネルギーを設計から減らす。
次にガードや保護装置等で接近と危害を防ぐ。それでも残るリスク、正しい使用・点検・訓練等の情報を提供する。
使用上の情報を落として「二つだけの方策」としない。
警告を増やすことだけで本質的な安全設計の検討を代替しない。
機械の製造側と使用側は各段階で役割を持ち、残留リスクがあるから全て使用者の責任という結論にはならない。

システム安全はハードウェアだけでなく、環境・人的条件を含める。
機能安全は必要な安全機能が正しく働くことによる安全で、通常の制御機能全てを意味しない。
安全工学は使用者と作業者双方の危害を扱い、レジリエンスは予防だけでなく
変化への対応や回復できる能力へ目を向ける。原典の対比表現を、全規格の排他的な定義へ広げない。

\section{リコール：危害防止の是正と賠償の成立を分ける}
\label{theory:safety:recall}
安全を損なう製品を把握したら、対象範囲と危害を調べ、回収、修理、交換、使用停止の連絡等を検討する。
劣化や表示の誤りでも是正が必要になり得る。
医薬品に別の品質管理法があることから、回収がないとは言えない。
PMDAは医薬品の回収情報を公表している。\cite{saf53_medicine_recall}

消安法の対象範囲、薬事等の別制度、行政への報告、購入者への連絡、社告を別の問いへ分ける。
対象が全て追跡されて回収できる場合の社告の扱いは、当時の手引きの細条件を要確認とする。
2010手引きの直接全文取得は成功しておらず、平成25年の条件を現行手引きだけから確定しない。\cite{saf53_recall2010}
回収した事実はPL責任の成立を直ちに示さない。欠陥・損害・因果関係等は別に検討する。
一方、賠償が成立しない可能性を、危害防止へ取り組まない根拠にしない。\cite{pl_statute,pl_qa}

\section{重大製品事故：報告と公表の起点を区別する}
\label{theory:safety:productaccident}
重大事故の代表類型は、死亡、治療に要する期間30日以上の重傷病、後遺障害、CO中毒、火災である。
製造・輸入側は事故の発生を知った段階から所定の期限内に報告する。
原因分析が未完という理由で期限を超えてよいとはしない。\cite{saf53_products2012}
事故の日、事業者が知った日、行政が報告を受理した日は一致するとは限らない。
後年取得した事業者用手引きは、知った日を含む10日を示すが、2018施行府令全文の照合と同一視しない。\cite{saf53_products2023}

行政の公表は事業者の報告と別の段階である。
2013年のフローでは、ガス・石油機器や製品起因が疑われる事故の迅速な公表と、
他の段階の報告受理後原則1週間という扱いを区別する。
「原則」を全案件に一律の法定絶対期限へ変えない。\cite{saf53_productflow2013}
通常有すべき安全性と、あらゆる条件での絶対安全は異なる。
合理的に予見できる使用を検討し、残るリスクを説明して減らすことが課題である。

長期使用の点検制度は経年劣化への対応を支える。
平成29年度の9品目は当時の指定として保持し、2021年等の対象見直しを遡及しない。
指定から外れたことも、劣化や保守の問題が消えるという意味ではない。\cite{saf53_oldnine}

\section{食品リスク：評価・管理・対話をつなぐ}
\label{theory:safety:foodroles}
科学的な健康影響を評価する機能と、許容する条件や施策を決め実施する管理の機能を分ける。
食品安全委員会の旧説明は、科学的な評価と管理行政の独立性を示す。\cite{saf53_food_roles2007}
評価結果だけから全ての政策が一意に決まるとは限らず、実行可能性、費用、便益や他の影響も検討する。
管理者の都合で評価結果を改変することも許されない。

リスクコミュニケーションでは、行政が情報を整理し対話の場を設け、
事業者は生産・流通・販売の情報を共有し、消費者は必要な情報を求め意見を述べ、
科学者は根拠と限界を分かりやすく説明する。
全過程における意見と情報の交換であって、結論を一方向に伝えて納得させるだけではない。\cite{saf_fsc_psych,saf_fsc_glossary}
リスク評価と管理の担当を入れ替えた肢は、この機能分担から判断できる。

\section{品質不正：記録の真正性と構造安全を別に検証する}
\label{theory:safety:quality}
測定記録を流用すれば、設計条件どおり施工したかを第三者が検証する根拠を失う。
記録不正だけで全ての構造物が倒壊すると断定できないが、実際の安全性の迅速な確認と
不正行為・管理過程の是正はともに必要である。
2015年の基礎ぐい報告は、責任体制と支持層到達等の技術判断を扱う。\cite{saf53_pile2015}

元請は施工全体を統括し、専門工事の担当は現場で得る情報から技術的に判断する。
設計図があるから現地の地盤確認を省略するとはしない。
機器が不調になったら、作業の停止や再取得、異常の報告と判断の経路を用意する。
ハードウェアの高度化だけでは、記録を偽る組織の圧力や曖昧な責任を取り除けない。
品質不正は他社からも教訓を得られる。トップの優先順位、部門間の役割、検査の独立性まで調べ、
個人の行為と組織条件の両方へ対応する。\cite{saf53_quality2020}

\section{作業の変更と軽微な危害の例外}
\label{theory:safety:workplace}
材料、設備、作業手順や人員の条件を変更したら、以前の評価の前提を確認する。
2006年の指針は、建設物や設備・材料・手順の変更時に調査し、疲労等の付加的な影響も考慮する。
対象の選定では、明らかに軽微な負傷・疾病しか予想されない平坦な歩行等を除外できる。
解説は軽微を医師の治療を要しない程度とし、過去に偶然軽い結果しか出なかったという条件とは区別する。\cite{saf53_workplace2006}

対策の優先順位と実行可能性を検討し、低減を続ける。
ALARPは、追加の低減による便益に対して費用等が著しく不均衡となるかを検討する原則で、
単なる安い対策の選択ではない。法令の最低基準を無視する理由にもならない。\cite{saf_parliament2006}

\section{ヒューマンエラー：結果の手前にある条件}
\label{theory:safety:human}
不適切な操作を認めても、「不注意」で説明を終えない。
知識・経験、注意、伝達、組織、手順の省略、状況への反射、混乱、錯覚、
心身の能力、疲労、単調さなどから、なぜその操作が起きたかを調べる。
教育資料の12分類はこれらの発生要因を含む。\cite{saf53_human12}
疲労や年齢による機能変化は、誤操作という行為自体と同じ概念ではない。
令和2年の問いは発生要因の分類に対応するため0項目を選ぶが、前文の9割以上という数値は
母集団・集計法が未確認である。全分野の普遍的な事故率として使わない。

対策は作業者への注意だけに限らない。認知しやすい表示、誤操作を防ぐ形状、
異常時の停止、負荷と休息、明確な連絡経路、責任分担と記録を組み合わせる。
原因の分類名を覚えるだけで、具体的な工程や条件を調べる課題を代替しない。

\section{2021行動計画：章と七つの取組を読む}
\label{theory:safety:maintenance2021}
第2次インフラ長寿命化計画は2021年6月18日に公表され、令和3〜7年度を計画期間とする。
必要施策の方向性にある七項目は、個別施設計画、点検診断・修繕更新、予算管理、体制、
新技術、情報基盤、基準類である。現状課題の別章と見出し数を混同しない。\cite{saf53_infra2021}
技術継承が独立した七項目の見出しにないことは、人材の確保や能力の継承が不要という意味ではない。
計画の階層を読み、見出しへの所属と実際の施策内容を区別する。
